\section{A dihedral CM host after real multiquadratic enlargement}
\label{sec:cubic-dihedral-cm-host}

\paragraph{Recovery status.}
This is a new recovery revision, written after the executor returned on
22 November 2026.  The earlier AC-G73-host manuscript was not present in
the restored checkpoint or the identified private stores.  The proof
below is reconstructed from the restored AC-G73 and the explicitly
retained algebraic argument; it is not represented as the missing
original bytes or as an unchanged historical freeze.

\begin{proposition}[AC-G73-host: commutator, cyclotomic exclusion, and restriction]
\label{AC-G73-host}
Let \(p>3\) be prime, let \(s\geq1\), and let \(T/\mathbb Q\) be a
totally real Galois extension with group
\[
 \operatorname{Gal}(T/\mathbb Q)
       =P\rtimes\langle c\rangle,\qquad
 P\simeq C_{p^s},\quad c^2=1,\quad crc^{-1}=r^{-1}.
 \tag{G73H.1}
\]
Write \(E=T^P\) for its quadratic subfield.  Let \(M/\mathbb Q\)
be any finite real multiquadratic extension, allowing \(M=\mathbb Q\),
and put
\[
                  K=\mathbb Q(\sqrt{-3}),\qquad H=KTM.
\]
Then \(H\) is CM with maximal real subfield \(TM\), and
\[
 \begin{split}
 [H:K]&=\frac{2p^s[M:\mathbb Q]}{[T\cap M:\mathbb Q]},\\
 T\cap M&\in\{\mathbb Q,E\},\\
 [\operatorname{Gal}(H/\mathbb Q),
       \operatorname{Gal}(H/\mathbb Q)]&=P,\\
                   H^{P}&=KEM .
 \end{split}                                                   \tag{G73H.2}
\]
Here \(P\) is identified with its natural subgroup in the compositum.
The last field is the maximal abelian subfield of \(H/\mathbb Q\);
its Galois group has exponent at most two.  Consequently
\[
                              \mu_p\not\subset H .              \tag{G73H.3}
\]
Every absolutely irreducible three-dimensional characteristic-zero
representation of \(G_K\) remains absolutely irreducible on \(G_H\).
Every character of exact order three of \(G_K\) retains that order
on \(G_H\).

If the actual inertia of such a representation at a prime \(q\)
is killed by \(KT/K\) as in AC-G73, it remains killed after replacing
\(KT\) by \(H\).  This assertion uses that actual inertia hypothesis;
it does not produce it.

Finally, put \(H'=H(\sqrt2)\).  Then all the preceding conclusions
remain true, and, with \(d=[H':H]\in\{1,2\}\),
\[
 [H':\mathbb Q]=d[H:\mathbb Q],\qquad
 [(H')^+:\mathbb Q]=d[H^+:\mathbb Q],\qquad
 \operatorname{rd}(H')\leq\sqrt8\,\operatorname{rd}(H).
 \tag{G73H.4}
\]
When \(d=1\), the root discriminant is unchanged.
\end{proposition}

\begin{proof}
The commutator subgroup of \(P\rtimes C_2\) is \(P\):
the quotient by \(P\) is abelian, and
\([c,r]=r^{-2}\) runs through \(P\) because \(p\) is odd.
Its maximal abelian subextension is therefore \(E\).
Since \(T\cap M\) is abelian over \(\mathbb Q\), it is contained
in \(E\); this proves the second assertion of (G73H.2).

Write \(J=T\cap M\).  The exact Galois group of \(TM/\mathbb Q\) is
the fiber product
\[
 \mathcal G=
 \{(g,a)\in\operatorname{Gal}(T/\mathbb Q)
                 \times\operatorname{Gal}(M/\mathbb Q):
                      g|_J=a|_J\}.
 \tag{G73H.5}
\]
This gives the degree formula for \(TM\) and embeds \(P\) as
\(\{(r,1):r\in P\}\).  The quotient \(\mathcal G/P\) is a
subgroup of the product of two elementary abelian two-groups,
so it is elementary abelian.  In particular
\([\mathcal G,\mathcal G]\subset P\).

There is an element of \(\mathcal G\) whose first coordinate is
the reflection \(c\).  If \(J=\mathbb Q\), take \((c,1)\).
If \(J=E\), choose any \(a\in\operatorname{Gal}(M/\mathbb Q)\)
restricting to the nontrivial automorphism of \(E\), and take
\((c,a)\).  Its conjugation on \(P\) is inversion.  Hence its
commutators with \((r,1)\) give all \(r^{-2}\), proving
\[
                         [\mathcal G,\mathcal G]=P.
 \tag{G73H.6}
\]
The fixed field of this subgroup is exactly \(EM\), by the
fiber-product description.

The field \(TM\) is totally real, whereas \(K\) is imaginary
quadratic; their intersection is \(\mathbb Q\).
Thus \(H/\mathbb Q\) has Galois group
\(C_2\times\mathcal G\).  Complex conjugation is its first
factor, and its fixed field is \(TM\), proving the CM assertion.
Equations (G73H.2) and the maximal-abelian-field assertion now
follow from (G73H.6).
If \(H\) contained a primitive \(p\)-th root of unity, restriction
would give a cyclic quotient of \(\operatorname{Gal}(H/\mathbb Q)\)
of order \(p-1>2\).  Every abelian quotient has exponent at most
two, a contradiction.  This proves (G73H.3).

For restriction of a three-dimensional representation, first
adjoin the finitely many real quadratic factors of \(M\) and
then \(E\), omitting any already present.  These are successive
cyclic extensions of degree two.  The last extension from
\(KME\) to \(KMT=H\) is cyclic of \(p\)-power degree.
Indeed \(T/E\) is cyclic of degree \(p^s\); intersection with
the multiquadratic field \(KME\) can only be \(E\).

We recall the exact representation argument used in AC-G73.
Restriction of an irreducible finite-dimensional representation
to a normal subgroup of finite index is semisimple: conjugates
of a simple submodule span the original representation.
For a cyclic quotient of order \(m\), over an algebraically
closed characteristic-zero coefficient field,
\[
 \operatorname{End}_{G'}(V)
   \simeq
   \bigoplus_{\nu\in\widehat{G/G'}}
             \operatorname{Hom}_{G}(V,V\otimes\nu).
 \tag{G73H.7}
\]
An isomorphism \(V\simeq V\otimes\nu\) forces
\(\nu^3=1\) by taking determinants.  If \(m=2\) or \(m\) is a
\(p\)-power with \(p>3\), only the trivial character can occur.
The commutant of the semisimple restriction is therefore
one-dimensional, and the restriction is irreducible.
Apply this successively in the displayed tower.
The argument after scalar extension proves absolute
irreducibility for any characteristic-zero coefficient field.

The degree \([H:K]\) has only the prime factors \(2\) and \(p\).
A cyclic cubic extension of \(K\) therefore has trivial
intersection with \(H\), proving persistence of an order-three
character.  No assertion is made here that an arbitrary
quadratic residual character survives an arbitrary choice of
\(M\); that requires the separate freshness arguments used in
AC-G81 and AC-G83.
An inertia action already trivial on \(G_{KT}\) is trivially
still trivial on its subgroup \(G_H\).

For the last assertion replace \(M\) by \(M(\sqrt2)\) throughout.
This remains real multiquadratic.  The degree and real-degree
identities in (G73H.4) are immediate.
If \(d=2\), the order
\(\mathcal O_H[\sqrt2]\) in \(H'\) has relative discriminant
generated by \(8\).  The relative field discriminant divides
this order discriminant, and the discriminant tower formula gives
\[
 |D_{H'}|\leq |D_H|^2\,8^{[H:\mathbb Q]}.
\]
Taking the \(1/(2[H:\mathbb Q])\)-th power proves (G73H.4).
If \(d=1\), nothing changes.
\end{proof}

\paragraph{Scope and geometric dictionary.}
AC-G73 supplies examples of \(T\) with a prescribed finite tame
inertia kernel and its own degree and discriminant bounds.
The present proposition is an exact field-theoretic refinement
after real multiquadratic enlargement.  It does not supply an
automorphic representation, a generic ordinary point, every
source level or frame condition, numerical counting bounds for
an arbitrary \(M\), or the actual characteristic-\(p\) Hecke-curve
containment.  The elementary degree identities alone do not
verify those numerical hypotheses or their parity conventions.

The finite Galois composita and restriction of representations
give the usual finite-cover/local-system interpretation in the
already audited SGA~1, Expos\'e~V, \S\S4 and~7 scope.
No new theorem is attributed to SGA.
The new proof uses the restored AC-G73 and elementary
Galois, representation, and discriminant arguments; it does not
claim a fresh independent proof of the class-field inputs
retained inside AC-G73.
